\section{Background}

\subsection{Multi-Task Post-Training as Task Allocation}
\label{sec:bg_allocation}

We post-train a policy $\pi$ on $K\ge 2$ tasks for a budget of $T$ rounds. Task $k$ has input distribution $\mathcal{X}_k$ and reward $r_k$, with rewards sharing a common scale. Let $J_k(\pi)$ denote the expected return of policy $\pi$ on task $k$, and let
$J(\pi)\coloneqq\sum_{k=1}^K J_k(\pi)$.
At round $t$, the learner splits a fixed-size mini-batch according to an \textit{allocation}
\[
p_t\in\Delta^K
\coloneqq
\left\{
p\in\mathbb{R}_+^K:
\sum_{k=1}^K p_k=1
\right\},
\]
and takes one update of the underlying RL algorithm, which is held fixed. Thus, $p_1,\ldots,p_T$ are the only control variables.

The batch under a fixed allocation is random, so we reason about its conditional mean effect on policy performance. Let $\pi_{t+1}(p)$ denote the policy obtained by applying the round-$t$ update to $\pi_t$ under allocation $p$, and let $\pi_{t+1}\coloneqq\pi_{t+1}(p_t)$. We define the \textit{task-wise one-step policy improvement}
\[
U_{t,k}(p)
\coloneqq
J_k(\pi_{t+1}(p))-J_k(\pi_t),
\]
and the aggregate improvement
$U_t(p)\coloneqq\sum_{k=1}^K U_{t,k}(p)$.
Realized improvements telescope across training:
\begin{equation}
\label{eq:bg_telescope}
J(\pi_{T+1})-J(\pi_1)
=
\sum_{t=1}^{T} U_t(p_t).
\end{equation}
Maximizing $\sum_{t=1}^T U_t(p_t)$ is therefore the finite-budget policy-improvement objective for multi-task post-training. Since $\pi_t$ changes throughout training, the response functions $U_t$ are generally non-stationary and are revealed only through policy updates.


\subsection{Task Allocation as Constrained Optimization}
\label{sec:bg_constrained_allocation}

A natural approach to task allocation is to optimize the mixture according to a model of how allocation affects subsequent training progress. Aioli and its predecessors ~\citep{chen2025aioli,ye2025datamixinglawsoptimizing,chen2023skillitdatadrivenskillsframework,xie2023doremioptimizingdatamixtures,fan2024dogedomainreweightinggeneralization} formalizes this idea through the following fomalization.
\begin{equation}
\label{eq:bg_lmo}
\max_{p_1,\ldots,p_T\in\Delta^K}\ \sum_{t=1}^{T} U_t(p_t)
\quad \text{s.t.}\quad
U_{t,k}(p_t)
    =
    \phi_k \left( \sum_{j=1}^{K} A_{t,kj}p_{t,j} \right)
\ \forall k,t.
\end{equation}
where $\phi_k$ specifies the shape of the scaling law. $A_t\in\mathbb{R}^{K\times K}$ models cross-task interactions. Through \Cref{eq:bg_lmo}, the current mixture is mapped to future task improvements through a parametric cross-task response matrix, and the next allocation is chosen according to these predicted gains. Aioli's allocation score for task $j$ aggregates its predicted effect across all tasks,
$s_{t,j}=\sum_{k=1}^{K}A_{t,kj}$, leading to an exponential update of the form
\begin{equation}
\label{eq:bg_aioli_update}
    p_{t+1,j}
    =
    \frac{
        p_{t,j}\exp\!\left(\eta s_{t,j}\right)
    }{
        \sum_{\ell=1}^{K}
        p_{t,\ell}\exp\!\left(\eta s_{t,\ell}\right)
    }.
\end{equation}

This formulation makes dynamic allocation tractable by prescribing a functional relationship between allocation and future training progress. \textbf{PDCM retains the constrained-optimization view, but does not require such a parametric response law.} Instead, it uses a single-step directional concavity condition to construct first-order constraints directly from the observed task-improvement.



% \subsection{Online Optimization with Long-Term Constraints}
% \label{sec:bg_online}

% This is \emph{online optimization with long-term constraints}: the learner commits to $p_t$ before the round's outcome, then sees first-order feedback---$U_t(p_t)$ and a supergradient of each $U_{t,k}$---from which one constraint $h_{t,k}:\DeltaK\to\R$ per task is built. Feasibility holds in \emph{aggregate}, not per round ($[K]\coloneqq\{1,\ldots,K\}$):
% \begin{equation} \label{eq:bg_oco}
% \text{maximize}\ \ \sum_{t=1}^{T} U_t(p_t)
% \quad\text{subject to}\quad
% \sum_{t=1}^{T} h_{t,k}(p_t)\le 0,
% \ \ \forall k\in[K].
% \end{equation}
% It is solved when both the regret against the best fixed allocation feasible for every revealed constraint and each cumulative violation $\pos{\sum_{t}h_{t,k}(p_t)}$ are sublinear in $T$. For concave $U_t$ and convex $h_{t,k}$ the standard solution is mirror descent~\citep{Beck2003MirrorDA} in a primal--dual loop~\citep{mahdavi2011trading,jenatton2015adaptivealgorithmsonlineconvex,wei2019onlineprimaldualmirrordescent}: $\lambda_{t,k}\ge0$ prices constraint $k$, the primal iterate takes an entropic mirror-\emph{ascent} step on $U_t(p)-\sum_{k}\lambda_{t,k}h_{t,k}(p)$, and $\lambda_{t,k}$ ascends on $h_{t,k}(p_t)$. Steps of order $T^{-1/2}$ give $O(\sqrt{T})$ regret with $O(T^{3/4})$ violation~\citep{mahdavi2011trading,jenatton2015adaptivealgorithmsonlineconvex}, both $O(\sqrt{T})$ once the multipliers stay bounded, as under a Slater condition~\citep{wei2019onlineprimaldualmirrordescent}.
